% ECONOMICS_PROBLEM: {"id": "H87-133", "title": "Global minimum taxes and real investment", "jel_code": "H87", "source_id": "OP-0475"}
% !TeX program = xelatex
\documentclass[11pt,a4paper]{article}
\usepackage[margin=23mm]{geometry}
\usepackage{fontspec}
\setmainfont{Latin Modern Roman}
\usepackage{amsmath,amssymb,mathtools}
\usepackage{xcolor,enumitem,booktabs,longtable,array,xurl,hyperref}
\definecolor{muted}{HTML}{53636C}
\hypersetup{colorlinks=true,urlcolor=blue,linkcolor=blue}
\setlength{\parindent}{0pt}
\setlength{\parskip}{5pt}
\newcommand{\R}{\mathbb R}
\newcommand{\E}{\mathbb E}
\newcommand{\N}{\mathbb N}
\newcommand{\ind}{\mathbf 1}
\newcommand{\Var}{\operatorname{Var}}
\newcommand{\Cov}{\operatorname{Cov}}
\newcommand{\supp}{\operatorname{supp}}
\DeclareMathOperator*{\argmin}{arg\,min}
\DeclareMathOperator*{\argmax}{arg\,max}
\newcommand{\entrymeta}[3]{{\sffamily\small\color{muted}\textbf{\texttt{#1}}\quad|\quad #2\par #3\par}}
\newcommand{\Problemref}[1]{\texttt{#1}}
\newcommand{\JELref}[2]{\texttt{#2} (JEL \texttt{#1})}
\begin{document}
\section*{Global minimum taxes and real investment}
\textbf{Problem ID:} \texttt{H87-133}\quad \textbf{JEL:} \texttt{H87}\par
% BEGIN ECONOMICS STATEMENT
\label{problem:OP-0475}
{\sffamily\small\color{muted}Original subject group: Taxation and social insurance\par}
\label{definitions:problem:30007757}
\textbf{Audit classification:} Background; proposed extension. All four authors and the July 31, 2026 author-hosted full text are verified. The NBER page/PDF failed to open, but a primary author PDF provides the model. Its fixed-policy quantitative answers do not establish this strategic extension as a published open problem.
\subsection*{Definitions and precise research target}
Fix \(J\geq2\) jurisdictions, firm types \(f=1,\ldots,F_0\) with specified masses \(n_f>0\) and per-firm outcomes, and a finite horizon \(H\). A tax regime \(m\) contains a declared minimum rate, scope, accounting base, priority of top-up claims, carve-outs and enforcement; legal terminology alone is not a model. Each government chooses an instrument vector \(t_j\) from compact \(\mathcal T_j\). Firms choose bounded tangible investment \(K_{fj}\), research or intangible investment \(D_{fj}\), location and reported profits \(\Pi_{fj}\), subject to declared production, resource and accounting constraints and real costs of shifting. Specify ownership and the government objective \(W_j\), including revenue use and foreign profit remittances.

For every policy profile \(t\), define the nonempty firm/market equilibrium correspondence \(\mathcal F(m,t)\), with a public selection rule \(s_m(t)\) or a declared worst-equilibrium payoff. This selection must be defined at off-equilibrium tax profiles as well. A government-policy equilibrium is a profile \(t^*\) such that
\[
 W_j(m,t^*,s_m(t^*))\geq
 W_j(m,(t_j,t^*_{-j}),s_m(t_j,t^*_{-j}))
 \quad\forall t_j\in\mathcal T_j,\ \forall j.
\]
The resulting equilibrium set \(\mathcal E(m)\) need not be nonempty merely because instrument sets are compact. Restrict comparison to regimes with nonempty sets and bounded outcomes. Define \(Z(e)=((\sum_fn_fK_{fj},\sum_fn_fD_{fj},\sum_fn_f\Pi_{fj},R_j,W_j))_{j=1}^J\), consistently discounted where appropriate. Relative to no-minimum regime \(m_0\), the set-valued counterfactual is
\[
 \mathcal D(m,m_0)=\{Z(e)-Z(e_0):e\in\mathcal E(m),\ e_0\in\mathcal E(m_0)\}.
\]
Selected-regime differences require a further public cross-regime selection rule. Characterize conditions making the aggregate investment coordinates positive or negative, and identify primitives or sharp bounds from firm/tax observations. Separate fixed-government-rate comparisons from strategic government adjustment. The cited paper already solves a particular quantitative reform model; the strategic government extension and uniform sign questions here are editorial.
\subsection*{Variants and assumption changes}
\begin{enumerate}
\item \textbf{Fixed tax rates versus tax competition (editorial).} Fix \(t=t^0\) and solve only firm behavior, then compare with the government equilibrium above. The latter is expressly a game-theoretic extension.
\item \textbf{Profit shifting versus intangible relocation (editorial).} Permit only accounting reassignment in one model; permit ownership and research location to change in another. These different constraints can change real-investment signs.
\item \textbf{Common selection versus all equilibria (editorial).} A selected \(\Delta Z\) describes one comparison; \(\mathcal D\) tests whether its sign survives equilibrium multiplicity.
\end{enumerate}
\subsection*{References and source audit}
Dyrda--Hong--Sajid--Steinberg, author-hosted July 31, 2026 draft: cover, model sections and reform exercises verified. The NBER identifier is independently indexed; direct NBER retrieval was blocked. Use the author PDF for content; no exact open-status claim is made.
\begin{enumerate}\raggedright
\item\hypertarget{definitions:source:30007757:1}{} Sebastian Dyrda; Guangbin Hong; Muhammad Ali Sajid; Joseph B. Steinberg. 2026. \emph{Global Ripple Effects of Corporate Tax Reforms}. Working Paper 34627; bibliographic record and abstract.
\url{https://www.nber.org/papers/w34627}.
DOI: \url{https://doi.org/10.3386/w34627}.
\item\hypertarget{definitions:source:30007757:2}{} Sebastian Dyrda; Guangbin Hong; Muhammad Ali Sajid; Joseph B. Steinberg. 2026. \emph{Global Ripple Effects of Corporate Tax Reforms}. Author draft dated July31,2026; cover, model and reform exercises.
\url{https://dyrda.info/files/DHSS_Global.pdf}.
\end{enumerate}

\subsection*{Common conventions: Source mathematical conventions}

These conventions apply to each entry unless explicitly replaced.

\textbf{Models and identification.} A primitive model \(m\) specifies all probability laws, feasible choices, information, equilibrium and measurement rules used by its target. The admissible class \(\mathcal M\) is fixed independently of outcomes. If \(P_O(m)\) is its observable probability law and \(\tau(m)\) the target, its identified set at observable law \(P\) is
\[
 \mathcal I_\tau(P)=\{\tau(m):m\in\mathcal M,\ P_O(m)=P\}.
\]
Point identification means that this set is a singleton. An empty set signals incompatibility of data and assumptions. Coordinate bounds need not describe the joint set. Bounds are sharp only if every retained target is attainable or arbitrarily approachable by compatible models. Finite bounds require explicit restrictions on unobserved quantities; unrestricted confounding, hidden wealth, extrapolation or arbitrary equilibrium selection can yield uninformative sets.

\textbf{Causal interventions.} A potential outcome \(Y_i(p)\) is the outcome under a complete policy regime \(p\), including timing, financing, information and market spillovers. The target population and units are fixed before treatment. With interference, use the complete assignment vector or a justified exposure mapping; individual no-interference notation alone cannot identify an economy-wide intervention. A difference of mean potential outcomes does not require the cross-world joint distribution. Conditioning on post-treatment migration, survival, enrollment or employment changes the estimand and needs separate justification.

\textbf{Statistical statements.} An estimator, confidence set, forecast or test specifies a sampling law, model class and loss/coverage criterion. Uniform \((1-\alpha)\)-coverage means \(\inf_{m\in\mathcal M}\Pr_m\{\tau(m)\in C_n\}\geq1-\alpha\), or an explicitly asymptotic version. The whole parameter space trivially has coverage, so informativeness requires a stated width, risk or power objective. A forecasting improvement is relative to a named benchmark, information set and evaluation population.

\textbf{Equilibrium and optimization.} An equilibrium comprises feasible plans, optimality, beliefs where needed, budget constraints and all market/resource clearing. The primitives or a stated benchmark define each correspondence \(\mathcal E(p;m)\). Nonemptiness is assumed or proved, never inferred just from compact policy sets. A selected-equilibrium target uses a declared selection rule; a robust target quantifies over all permitted equilibria. Use suprema or infima unless attainment is established. An \(\varepsilon\)-optimal policy is feasible and within \(\varepsilon\) of the finite optimum value. Report an empty feasible set separately.

\textbf{Welfare and accounting.} Normative utility, interpersonal normalization, social weights, numeraire, discounting and the use of public receipts are primitives. Behavioral utility need not coincide with the welfare criterion. Household welfare includes ownership income and rents once; adding profits again to compensating variation generally double-counts them. A monetary equivalent is defined by an explicit transfer and price convention, with monotonicity and range conditions. Average effects, mechanism contributions, transfers and welfare losses are different targets. Infinite sums require convergence and dynamic feasible plans require terminal or no-Ponzi conditions.

\textbf{Source versions.} A working-paper date, revision date and journal publication year are distinguished. A DOI for a journal version is not automatically the DOI of an earlier manuscript. Locators refer to the stated version and printed pages where specified. A metadata-only check does not verify a claim about a full-text conclusion. Every entry's source subsection narrows the provenance claim accordingly.
% END ECONOMICS STATEMENT
\end{document}
